% ---------------------------------------------------------------------------
\section{Executive Summary}
% ---------------------------------------------------------------------------

This project delivers a valuation capability for a residential listing platform in
Ames, Iowa. At the moment a property is listed, the platform must display a price
and indicate how much confidence an agent or buyer should place in it. The answer is
produced by a regularised regression model on the logarithm of sale price, wrapped in
a calibrated prediction interval, exposed through an application programming
interface (API) and a browser console, and surrounded by a monitoring layer that
compares incoming listing populations with the training population and tracks
realised error once sale prices arrive.



Table~\ref{tab:kpis} in Section~\ref{sec:business} sets out the KPI framework with
targets and achieved values; the headline numbers are repeated here. The selected
model explains 93.7 per cent of out-of-fold price variance and prices a
typical property within 13,537 \USD{} of its realised value. Nearly 70 per cent of
estimates fall inside the 15,000 \USD{} tolerance band that the business case treats
as the threshold for an agent-ready estimate, and the conformal intervals are
calibrated almost exactly at their nominal levels. Performance is stable when the
model is refitted on 2006--2009 sales and applied to 2010 sales, the closest
available approximation to genuine forward deployment.

Three findings shape how the product should be used. First, the simplest regularised
linear model outperformed both gradient-boosted trees and a random forest: with 1,458
transactions and a dense set of well-measured attributes, the variance cost of
flexible learners exceeds the bias they remove \citep{hastie2009esl}. Second,
accuracy and calibration are separate objectives. The model remains accurate
throughout the monitoring replay, yet its 90 per cent interval covers only 79.7 per
cent of sales in the final window, because the macro-financial context moves outside
the reference sample; a point-estimate dashboard would not have detected this.
Third, error is not uniform: the luxury band above 300,000 \USD{} is under-predicted
by 3.9 per cent with a mean absolute error of 31,238 \USD, more than three times that
of the entry band.

The deployment is verified locally rather than publicly. The API, console and
monitoring report run end to end against the trained artefacts and are containerised,
and the Render and Hugging Face configurations are complete. No cloud credentials
were available, so no public URL is claimed anywhere in this report. This is a
delivery constraint, not an analytical one, and Section~\ref{sec:deployment} states
exactly what remains for the account owner.
